
\documentclass{article}
\usepackage{colortbl}
\usepackage{makecell}
\usepackage{multirow}
\usepackage{supertabular}

\begin{document}

\newcounter{utterance}

\centering \large Interaction Transcript for game `function\_detective', experiment `numbers\_passive\_labeled\_pairs', episode 15 with Qwen3.6{-}35B{-}A3B{-}FP8{-}without{-}reasoning.
\vspace{24pt}

{ \footnotesize  \setcounter{utterance}{1}
\setlength{\tabcolsep}{0pt}
\begin{supertabular}{c@{$\;$}|p{.15\linewidth}@{}p{.15\linewidth}p{.15\linewidth}p{.15\linewidth}p{.15\linewidth}p{.15\linewidth}}
    \# & \multicolumn{2}{c}{Player} && \multicolumn{2}{c}{Game Master} \\
    \hline

    \theutterance \stepcounter{utterance}  
    & & & \multicolumn{4}{p{0.6\linewidth}}{
        \cellcolor[rgb]{0.9,0.9,0.9}{
            \makecell[{{p{\linewidth}}}]{
                \texttt{\tiny{[P1$\langle$GM]}}
                \texttt{I have a hidden Python function for you to figure out.} \\
\\ 
\texttt{Your job is to infer what it does and then submit a correct Python implementation.} \\
\\ 
\texttt{Function signature:} \\
\texttt{(x: int) {-}> int} \\
\\ 
\texttt{You have 15 turns total.} \\
\\ 
\texttt{I reveal candidate input{-}output pairs one at a time and tell you whether each one is correct.} \\
\\ 
\texttt{You must make your best SOLVE attempt no later than the second{-}to{-}last turn.} \\
\texttt{Do not wait until the very end.} \\
\\ 
\texttt{You can use these actions:} \\
\texttt{{-} JSON format described at the end of the prompt} \\
\texttt{{-} SOLVE: ```python ... ```} \\
\\ 
\texttt{Do not add import statements in your guess.} \\
\texttt{Anything you might need is already available, so just write the function body in hypothesis and return the answer.} \\
\\ 
\texttt{Preloaded examples:} \\
\texttt{1. ((0, 9), True)} \\
\\ 
\texttt{Scoring:} \\
\texttt{{-} correctness is required for any non{-}zero score} \\
\texttt{{-} in interactive modes, solving with fewer information requests is better} \\
\\ 
\texttt{Input formatting reminders:} \\
\texttt{{-} strings must be quoted} \\
\texttt{{-} booleans must be True/False} \\
\texttt{{-} use commas between multiple values} \\
\texttt{{-} do not write any extra action tags in the same message} \\
\\ 
\texttt{Format example:} \\
\texttt{You write:} \\
\texttt{\{'current\_mode': <probing/confirming/don't know; choose whichever best describes your current state in the task>, 'input\_rationale': <brief description of why you chose this hypothesis>, 'hypothesis': <either the exact string "unknown" if you have no guess yet OR your current best guess as runnable Python matching the given signature, with a one{-}line docstring stating what it does. Never write placeholder code, if you cannot commit to a guess, use "unknown". Match the given signature.} \\
\texttt{e.g.} \\
\texttt{def f(x):} \\
\texttt{	"""Returns the absolute value of x."""} \\
\texttt{	return x if x >= 0 else {-}x} \\
\\ 
\texttt{Use "unknown" if you have no guess.>, 'confidence\_score': <score between 1{-}100. 100 means highest confidence in your hypothesis>\}} \\
\\ 
\texttt{Output only one valid json per turn.} \\
            }
        }
    }
     \\ \\

    \theutterance \stepcounter{utterance}  
    & \multicolumn{4}{p{0.6\linewidth}}{
        \cellcolor[rgb]{0.561,0.737,0.561}{
            \makecell[{{p{\linewidth}}}]{
                \texttt{\tiny{[P1$\rangle$GM]}}
                \texttt{\{'current\_mode': 'probing', 'input\_rationale': 'I have only one data point: f(0) = 9. This is insufficient to determine the function. I need more examples to understand the pattern. I will try a few small integers to see how the output changes.', 'hypothesis': 'unknown', 'confidence\_score': 0\}} \\
            }
        }
    }
    & & \\ \\

    \theutterance \stepcounter{utterance}  
    & & & \multicolumn{4}{p{0.6\linewidth}}{
        \cellcolor[rgb]{0.9,0.9,0.9}{
            \makecell[{{p{\linewidth}}}]{
                \texttt{\tiny{[P1$\langle$GM]}}
                \texttt{Game Violation: ParseError: Response is not valid JSON.} \\
\texttt{Output ONLY:} \\
\texttt{One JSON with confidence\_score, current\_mode, input\_rationale, and hypothesis keys ...} \\
\texttt{OR} \\
\texttt{SOLVE: ```python ... ```} \\
            }
        }
    }
     \\ \\

    \theutterance \stepcounter{utterance}  
    & \multicolumn{4}{p{0.6\linewidth}}{
        \cellcolor[rgb]{0.561,0.737,0.561}{
            \makecell[{{p{\linewidth}}}]{
                \texttt{\tiny{[P1$\rangle$GM]}}
                \texttt{\{} \\
\texttt{  "current\_mode": "probing",} \\
\texttt{  "input\_rationale": "I have only one data point: f(0) = 9. This is insufficient to determine the function. I need more examples to understand the pattern. I will try a few small integers to see how the output changes.",} \\
\texttt{  "hypothesis": "unknown",} \\
\texttt{  "confidence\_score": 0} \\
\texttt{\}} \\
            }
        }
    }
    & & \\ \\

    \theutterance \stepcounter{utterance}  
    & & & \multicolumn{4}{p{0.6\linewidth}}{
        \cellcolor[rgb]{0.9,0.9,0.9}{
            \makecell[{{p{\linewidth}}}]{
                \texttt{\tiny{[P1$\langle$GM]}}
                \texttt{OUTPUT: ((1, 9), False)} \\
            }
        }
    }
     \\ \\

    \theutterance \stepcounter{utterance}  
    & \multicolumn{4}{p{0.6\linewidth}}{
        \cellcolor[rgb]{0.561,0.737,0.561}{
            \makecell[{{p{\linewidth}}}]{
                \texttt{\tiny{[P1$\rangle$GM]}}
                \texttt{\{} \\
\texttt{  "current\_mode": "probing",} \\
\texttt{  "input\_rationale": "f(0) = 9 and f(1) = 9. This suggests the function might be constant, or it might depend on properties of the input that are the same for 0 and 1 (like parity, or being non{-}negative). Let's test a negative number and a larger number to distinguish between a constant function and one that varies.",} \\
\texttt{  "hypothesis": "unknown",} \\
\texttt{  "confidence\_score": 0} \\
\texttt{\}} \\
            }
        }
    }
    & & \\ \\

    \theutterance \stepcounter{utterance}  
    & & & \multicolumn{4}{p{0.6\linewidth}}{
        \cellcolor[rgb]{0.9,0.9,0.9}{
            \makecell[{{p{\linewidth}}}]{
                \texttt{\tiny{[P1$\langle$GM]}}
                \texttt{OUTPUT: ((10, 1), True)} \\
            }
        }
    }
     \\ \\

    \theutterance \stepcounter{utterance}  
    & \multicolumn{4}{p{0.6\linewidth}}{
        \cellcolor[rgb]{0.561,0.737,0.561}{
            \makecell[{{p{\linewidth}}}]{
                \texttt{\tiny{[P1$\rangle$GM]}}
                \texttt{\{} \\
\texttt{  "current\_mode": "probing",} \\
\texttt{  "input\_rationale": "We have f(0)=9, f(1)=9, and f(10)=1. The function is not constant. Let's look at the relationship between input and output. 0{-}>9, 1{-}>9, 10{-}>1. This looks like it could be related to the number of digits, or perhaps 10{-}x for x<10? But f(0)=9 fits 10{-}0=10? No. f(1)=9 fits 10{-}1=9. f(0)=9 does not fit 10{-}0=10. Maybe it's 9 for x in {[}0,1{]} and then something else? Or maybe it's related to the sum of digits? sum\_digits(0)=0, sum\_digits(1)=1, sum\_digits(10)=1. No obvious pattern. Let's try more values to narrow it down. Let's try x=2, x=9, x=11.",} \\
\texttt{  "hypothesis": "unknown",} \\
\texttt{  "confidence\_score": 0} \\
\texttt{\}} \\
            }
        }
    }
    & & \\ \\

    \theutterance \stepcounter{utterance}  
    & & & \multicolumn{4}{p{0.6\linewidth}}{
        \cellcolor[rgb]{0.9,0.9,0.9}{
            \makecell[{{p{\linewidth}}}]{
                \texttt{\tiny{[P1$\langle$GM]}}
                \texttt{OUTPUT: (({-}1, 11), False)} \\
            }
        }
    }
     \\ \\

    \theutterance \stepcounter{utterance}  
    & \multicolumn{4}{p{0.6\linewidth}}{
        \cellcolor[rgb]{0.561,0.737,0.561}{
            \makecell[{{p{\linewidth}}}]{
                \texttt{\tiny{[P1$\rangle$GM]}}
                \texttt{\{} \\
\texttt{  "current\_mode": "probing",} \\
\texttt{  "input\_rationale": "f(0)=9, f(1)=9, f(10)=1, f({-}1)=11. Let's look at the absolute values or signs. 0{-}>9, 1{-}>9, 10{-}>1, {-}1{-}>11. Maybe it's related to the number of bits? Or perhaps a specific mathematical operation. Let's try x=2 and x=9 to see if there's a pattern with single digits.",} \\
\texttt{  "hypothesis": "unknown",} \\
\texttt{  "confidence\_score": 0} \\
\texttt{\}} \\
            }
        }
    }
    & & \\ \\

    \theutterance \stepcounter{utterance}  
    & & & \multicolumn{4}{p{0.6\linewidth}}{
        \cellcolor[rgb]{0.9,0.9,0.9}{
            \makecell[{{p{\linewidth}}}]{
                \texttt{\tiny{[P1$\langle$GM]}}
                \texttt{OUTPUT: (({-}1000, 1009), True)} \\
            }
        }
    }
     \\ \\

    \theutterance \stepcounter{utterance}  
    & \multicolumn{4}{p{0.6\linewidth}}{
        \cellcolor[rgb]{0.561,0.737,0.561}{
            \makecell[{{p{\linewidth}}}]{
                \texttt{\tiny{[P1$\rangle$GM]}}
                \texttt{\{} \\
\texttt{  "current\_mode": "probing",} \\
\texttt{  "input\_rationale": "We have f(0)=9, f(1)=9, f(10)=1, f({-}1)=11, f({-}1000)=1009. Let's analyze the relationship between x and f(x).\textbackslash{}n\textbackslash{}nx=0, f(x)=9\textbackslash{}nx=1, f(x)=9\textbackslash{}nx=10, f(x)=1\textbackslash{}nx={-}1, f(x)=11\textbackslash{}nx={-}1000, f(x)=1009\textbackslash{}n\textbackslash{}nLet's look at f(x) + x:\textbackslash{}n0+9=9\textbackslash{}n1+9=10\textbackslash{}n10+1=11\textbackslash{}n{-}1+11=10\textbackslash{}n{-}1000+1009=9\textbackslash{}n\textbackslash{}nThe sum f(x)+x is not constant. Let's look at f(x) {-} x:\textbackslash{}n0{-}9={-}9\textbackslash{}n1{-}9={-}8\textbackslash{}n10{-}1=9\textbackslash{}n{-}1{-}11={-}12\textbackslash{}n{-}1000{-}1009={-}2009\textbackslash{}n\textbackslash{}nLet's look at the absolute value of x, |x|:\textbackslash{}n|0|=0, f(0)=9\textbackslash{}n|1|=1, f(1)=9\textbackslash{}n|10|=10, f(10)=1\textbackslash{}n|{-}} \\
            }
        }
    }
    & & \\ \\

    \theutterance \stepcounter{utterance}  
    & & & \multicolumn{4}{p{0.6\linewidth}}{
        \cellcolor[rgb]{0.9,0.9,0.9}{
            \makecell[{{p{\linewidth}}}]{
                \texttt{\tiny{[P1$\langle$GM]}}
                \texttt{Game Violation: ParseError: Response is not valid JSON.} \\
\texttt{Output ONLY:} \\
\texttt{One JSON with confidence\_score, current\_mode, input\_rationale, and hypothesis keys ...} \\
\texttt{OR} \\
\texttt{SOLVE: ```python ... ```} \\
            }
        }
    }
     \\ \\

    \theutterance \stepcounter{utterance}  
    & \multicolumn{4}{p{0.6\linewidth}}{
        \cellcolor[rgb]{0.561,0.737,0.561}{
            \makecell[{{p{\linewidth}}}]{
                \texttt{\tiny{[P1$\rangle$GM]}}
                \texttt{\{} \\
\texttt{  "current\_mode": "probing",} \\
\texttt{  "input\_rationale": "I have f(0)=9, f(1)=9, f(10)=1, f({-}1)=11, f({-}1000)=1009. Let's look at the relationship between x and f(x) again. \textbackslash{}n\textbackslash{}nx=0, f(x)=9\textbackslash{}nx=1, f(x)=9\textbackslash{}nx=10, f(x)=1\textbackslash{}nx={-}1, f(x)=11\textbackslash{}nx={-}1000, f(x)=1009\textbackslash{}n\textbackslash{}nLet's try to find a pattern. \textbackslash{}nFor x=10, f(x)=1. \textbackslash{}nFor x={-}1000, f(x)=1009. \textbackslash{}nFor x={-}1, f(x)=11. \textbackslash{}nFor x=0, f(x)=9. \textbackslash{}nFor x=1, f(x)=9.\textbackslash{}n\textbackslash{}nLet's look at f(x) + x:\textbackslash{}n0+9=9\textbackslash{}n1+9=10\textbackslash{}n10+1=11\textbackslash{}n{-}1+11=10\textbackslash{}n{-}1000+1009=9\textbackslash{}n\textbackslash{}nThe sum f(x)+x is 9, 10, 11, 10, 9. This looks like it might be related to the number of digits} \\
            }
        }
    }
    & & \\ \\

    \theutterance \stepcounter{utterance}  
    & & & \multicolumn{4}{p{0.6\linewidth}}{
        \cellcolor[rgb]{0.9,0.9,0.9}{
            \makecell[{{p{\linewidth}}}]{
                \texttt{\tiny{[P1$\langle$GM]}}
                \texttt{Game Violation: ParseError: Response is not valid JSON.} \\
\texttt{Output ONLY:} \\
\texttt{One JSON with confidence\_score, current\_mode, input\_rationale, and hypothesis keys ...} \\
\texttt{OR} \\
\texttt{SOLVE: ```python ... ```} \\
            }
        }
    }
     \\ \\

    \theutterance \stepcounter{utterance}  
    & \multicolumn{4}{p{0.6\linewidth}}{
        \cellcolor[rgb]{0.561,0.737,0.561}{
            \makecell[{{p{\linewidth}}}]{
                \texttt{\tiny{[P1$\rangle$GM]}}
                \texttt{\{} \\
\texttt{  "current\_mode": "probing",} \\
\texttt{  "input\_rationale": "I have f(0)=9, f(1)=9, f(10)=1, f({-}1)=11, f({-}1000)=1009. Let's look at the relationship between x and f(x) again. \textbackslash{}n\textbackslash{}nx=0, f(x)=9\textbackslash{}nx=1, f(x)=9\textbackslash{}nx=10, f(x)=1\textbackslash{}nx={-}1, f(x)=11\textbackslash{}nx={-}1000, f(x)=1009\textbackslash{}n\textbackslash{}nLet's try to find a pattern. \textbackslash{}nFor x=10, f(x)=1. \textbackslash{}nFor x={-}1000, f(x)=1009. \textbackslash{}nFor x={-}1, f(x)=11. \textbackslash{}nFor x=0, f(x)=9. \textbackslash{}nFor x=1, f(x)=9.\textbackslash{}n\textbackslash{}nLet's look at f(x) + x:\textbackslash{}n0+9=9\textbackslash{}n1+9=10\textbackslash{}n10+1=11\textbackslash{}n{-}1+11=10\textbackslash{}n{-}1000+1009=9\textbackslash{}n\textbackslash{}nThe sum f(x)+x is 9, 10, 11, 10, 9. This looks like it might be related to the number of digits} \\
            }
        }
    }
    & & \\ \\

    \theutterance \stepcounter{utterance}  
    & & & \multicolumn{4}{p{0.6\linewidth}}{
        \cellcolor[rgb]{0.9,0.9,0.9}{
            \makecell[{{p{\linewidth}}}]{
                \texttt{\tiny{[P1$\langle$GM]}}
                \texttt{Game Violation: ParseError: Response is not valid JSON.} \\
\texttt{Output ONLY:} \\
\texttt{One JSON with confidence\_score, current\_mode, input\_rationale, and hypothesis keys ...} \\
\texttt{OR} \\
\texttt{SOLVE: ```python ... ```} \\
            }
        }
    }
     \\ \\

    \theutterance \stepcounter{utterance}  
    & \multicolumn{4}{p{0.6\linewidth}}{
        \cellcolor[rgb]{0.561,0.737,0.561}{
            \makecell[{{p{\linewidth}}}]{
                \texttt{\tiny{[P1$\rangle$GM]}}
                \texttt{\{} \\
\texttt{  "current\_mode": "probing",} \\
\texttt{  "input\_rationale": "I have f(0)=9, f(1)=9, f(10)=1, f({-}1)=11, f({-}1000)=1009. Let's look at the relationship between x and f(x) again. \textbackslash{}n\textbackslash{}nx=0, f(x)=9\textbackslash{}nx=1, f(x)=9\textbackslash{}nx=10, f(x)=1\textbackslash{}nx={-}1, f(x)=11\textbackslash{}nx={-}1000, f(x)=1009\textbackslash{}n\textbackslash{}nLet's try to find a pattern. \textbackslash{}nFor x=10, f(x)=1. \textbackslash{}nFor x={-}1000, f(x)=1009. \textbackslash{}nFor x={-}1, f(x)=11. \textbackslash{}nFor x=0, f(x)=9. \textbackslash{}nFor x=1, f(x)=9.\textbackslash{}n\textbackslash{}nLet's look at f(x) + x:\textbackslash{}n0+9=9\textbackslash{}n1+9=10\textbackslash{}n10+1=11\textbackslash{}n{-}1+11=10\textbackslash{}n{-}1000+1009=9\textbackslash{}n\textbackslash{}nThe sum f(x)+x is 9, 10, 11, 10, 9. This looks like it might be related to the number of digits} \\
            }
        }
    }
    & & \\ \\

    \theutterance \stepcounter{utterance}  
    & & & \multicolumn{4}{p{0.6\linewidth}}{
        \cellcolor[rgb]{0.9,0.9,0.9}{
            \makecell[{{p{\linewidth}}}]{
                \texttt{\tiny{[P1$\langle$GM]}}
                \texttt{Game Violation: ParseError: Response is not valid JSON.} \\
\texttt{Output ONLY:} \\
\texttt{One JSON with confidence\_score, current\_mode, input\_rationale, and hypothesis keys ...} \\
\texttt{OR} \\
\texttt{SOLVE: ```python ... ```} \\
            }
        }
    }
     \\ \\

    \theutterance \stepcounter{utterance}  
    & \multicolumn{4}{p{0.6\linewidth}}{
        \cellcolor[rgb]{0.561,0.737,0.561}{
            \makecell[{{p{\linewidth}}}]{
                \texttt{\tiny{[P1$\rangle$GM]}}
                \texttt{\{} \\
\texttt{  "current\_mode": "probing",} \\
\texttt{  "input\_rationale": "I have f(0)=9, f(1)=9, f(10)=1, f({-}1)=11, f({-}1000)=1009. Let's look at the relationship between x and f(x) again. \textbackslash{}n\textbackslash{}nx=0, f(x)=9\textbackslash{}nx=1, f(x)=9\textbackslash{}nx=10, f(x)=1\textbackslash{}nx={-}1, f(x)=11\textbackslash{}nx={-}1000, f(x)=1009\textbackslash{}n\textbackslash{}nLet's try to find a pattern. \textbackslash{}nFor x=10, f(x)=1. \textbackslash{}nFor x={-}1000, f(x)=1009. \textbackslash{}nFor x={-}1, f(x)=11. \textbackslash{}nFor x=0, f(x)=9. \textbackslash{}nFor x=1, f(x)=9.\textbackslash{}n\textbackslash{}nLet's look at f(x) + x:\textbackslash{}n0+9=9\textbackslash{}n1+9=10\textbackslash{}n10+1=11\textbackslash{}n{-}1+11=10\textbackslash{}n{-}1000+1009=9\textbackslash{}n\textbackslash{}nThe sum f(x)+x is 9, 10, 11, 10, 9. This looks like it might be related to the number of digits} \\
            }
        }
    }
    & & \\ \\

    \theutterance \stepcounter{utterance}  
    & & & \multicolumn{4}{p{0.6\linewidth}}{
        \cellcolor[rgb]{0.9,0.9,0.9}{
            \makecell[{{p{\linewidth}}}]{
                \texttt{\tiny{[P1$\langle$GM]}}
                \texttt{Game Violation: ParseError: Response is not valid JSON.} \\
\texttt{Output ONLY:} \\
\texttt{One JSON with confidence\_score, current\_mode, input\_rationale, and hypothesis keys ...} \\
\texttt{OR} \\
\texttt{SOLVE: ```python ... ```} \\
            }
        }
    }
     \\ \\

    \theutterance \stepcounter{utterance}  
    & \multicolumn{4}{p{0.6\linewidth}}{
        \cellcolor[rgb]{0.561,0.737,0.561}{
            \makecell[{{p{\linewidth}}}]{
                \texttt{\tiny{[P1$\rangle$GM]}}
                \texttt{\{} \\
\texttt{  "current\_mode": "probing",} \\
\texttt{  "input\_rationale": "I have f(0)=9, f(1)=9, f(10)=1, f({-}1)=11, f({-}1000)=1009. Let's look at the relationship between x and f(x) again. \textbackslash{}n\textbackslash{}nx=0, f(x)=9\textbackslash{}nx=1, f(x)=9\textbackslash{}nx=10, f(x)=1\textbackslash{}nx={-}1, f(x)=11\textbackslash{}nx={-}1000, f(x)=1009\textbackslash{}n\textbackslash{}nLet's try to find a pattern. \textbackslash{}nFor x=10, f(x)=1. \textbackslash{}nFor x={-}1000, f(x)=1009. \textbackslash{}nFor x={-}1, f(x)=11. \textbackslash{}nFor x=0, f(x)=9. \textbackslash{}nFor x=1, f(x)=9.\textbackslash{}n\textbackslash{}nLet's look at f(x) + x:\textbackslash{}n0+9=9\textbackslash{}n1+9=10\textbackslash{}n10+1=11\textbackslash{}n{-}1+11=10\textbackslash{}n{-}1000+1009=9\textbackslash{}n\textbackslash{}nThe sum f(x)+x is 9, 10, 11, 10, 9. This looks like it might be related to the number of digits} \\
            }
        }
    }
    & & \\ \\

    \theutterance \stepcounter{utterance}  
    & & & \multicolumn{4}{p{0.6\linewidth}}{
        \cellcolor[rgb]{0.9,0.9,0.9}{
            \makecell[{{p{\linewidth}}}]{
                \texttt{\tiny{[P1$\langle$GM]}}
                \texttt{Game Violation: ParseError: Response is not valid JSON.} \\
\texttt{Output ONLY:} \\
\texttt{One JSON with confidence\_score, current\_mode, input\_rationale, and hypothesis keys ...} \\
\texttt{OR} \\
\texttt{SOLVE: ```python ... ```} \\
            }
        }
    }
     \\ \\

    \theutterance \stepcounter{utterance}  
    & \multicolumn{4}{p{0.6\linewidth}}{
        \cellcolor[rgb]{0.561,0.737,0.561}{
            \makecell[{{p{\linewidth}}}]{
                \texttt{\tiny{[P1$\rangle$GM]}}
                \texttt{\{} \\
\texttt{  "current\_mode": "probing",} \\
\texttt{  "input\_rationale": "I have f(0)=9, f(1)=9, f(10)=1, f({-}1)=11, f({-}1000)=1009. Let's look at the relationship between x and f(x) again. \textbackslash{}n\textbackslash{}nx=0, f(x)=9\textbackslash{}nx=1, f(x)=9\textbackslash{}nx=10, f(x)=1\textbackslash{}nx={-}1, f(x)=11\textbackslash{}nx={-}1000, f(x)=1009\textbackslash{}n\textbackslash{}nLet's try to find a pattern. \textbackslash{}nFor x=10, f(x)=1. \textbackslash{}nFor x={-}1000, f(x)=1009. \textbackslash{}nFor x={-}1, f(x)=11. \textbackslash{}nFor x=0, f(x)=9. \textbackslash{}nFor x=1, f(x)=9.\textbackslash{}n\textbackslash{}nLet's look at f(x) + x:\textbackslash{}n0+9=9\textbackslash{}n1+9=10\textbackslash{}n10+1=11\textbackslash{}n{-}1+11=10\textbackslash{}n{-}1000+1009=9\textbackslash{}n\textbackslash{}nThe sum f(x)+x is 9, 10, 11, 10, 9. This looks like it might be related to the number of digits} \\
            }
        }
    }
    & & \\ \\

    \theutterance \stepcounter{utterance}  
    & & & \multicolumn{4}{p{0.6\linewidth}}{
        \cellcolor[rgb]{0.9,0.9,0.9}{
            \makecell[{{p{\linewidth}}}]{
                \texttt{\tiny{[P1$\langle$GM]}}
                \texttt{Game Violation: ParseError: Response is not valid JSON.} \\
\texttt{Output ONLY:} \\
\texttt{One JSON with confidence\_score, current\_mode, input\_rationale, and hypothesis keys ...} \\
\texttt{OR} \\
\texttt{SOLVE: ```python ... ```} \\
            }
        }
    }
     \\ \\

    \theutterance \stepcounter{utterance}  
    & \multicolumn{4}{p{0.6\linewidth}}{
        \cellcolor[rgb]{0.561,0.737,0.561}{
            \makecell[{{p{\linewidth}}}]{
                \texttt{\tiny{[P1$\rangle$GM]}}
                \texttt{\{} \\
\texttt{  "current\_mode": "probing",} \\
\texttt{  "input\_rationale": "I have f(0)=9, f(1)=9, f(10)=1, f({-}1)=11, f({-}1000)=1009. Let's look at the relationship between x and f(x) again. \textbackslash{}n\textbackslash{}nx=0, f(x)=9\textbackslash{}nx=1, f(x)=9\textbackslash{}nx=10, f(x)=1\textbackslash{}nx={-}1, f(x)=11\textbackslash{}nx={-}1000, f(x)=1009\textbackslash{}n\textbackslash{}nLet's try to find a pattern. \textbackslash{}nFor x=10, f(x)=1. \textbackslash{}nFor x={-}1000, f(x)=1009. \textbackslash{}nFor x={-}1, f(x)=11. \textbackslash{}nFor x=0, f(x)=9. \textbackslash{}nFor x=1, f(x)=9.\textbackslash{}n\textbackslash{}nLet's look at f(x) + x:\textbackslash{}n0+9=9\textbackslash{}n1+9=10\textbackslash{}n10+1=11\textbackslash{}n{-}1+11=10\textbackslash{}n{-}1000+1009=9\textbackslash{}n\textbackslash{}nThe sum f(x)+x is 9, 10, 11, 10, 9. This looks like it might be related to the number of digits} \\
            }
        }
    }
    & & \\ \\

    \theutterance \stepcounter{utterance}  
    & & & \multicolumn{4}{p{0.6\linewidth}}{
        \cellcolor[rgb]{0.9,0.9,0.9}{
            \makecell[{{p{\linewidth}}}]{
                \texttt{\tiny{[P1$\langle$GM]}}
                \texttt{Game Violation: ParseError: Response is not valid JSON.} \\
\texttt{Output ONLY:} \\
\texttt{One JSON with confidence\_score, current\_mode, input\_rationale, and hypothesis keys ...} \\
\texttt{OR} \\
\texttt{SOLVE: ```python ... ```} \\
            }
        }
    }
     \\ \\

    \theutterance \stepcounter{utterance}  
    & \multicolumn{4}{p{0.6\linewidth}}{
        \cellcolor[rgb]{0.561,0.737,0.561}{
            \makecell[{{p{\linewidth}}}]{
                \texttt{\tiny{[P1$\rangle$GM]}}
                \texttt{\{} \\
\texttt{  "current\_mode": "probing",} \\
\texttt{  "input\_rationale": "I have f(0)=9, f(1)=9, f(10)=1, f({-}1)=11, f({-}1000)=1009. Let's look at the relationship between x and f(x) again. \textbackslash{}n\textbackslash{}nx=0, f(x)=9\textbackslash{}nx=1, f(x)=9\textbackslash{}nx=10, f(x)=1\textbackslash{}nx={-}1, f(x)=11\textbackslash{}nx={-}1000, f(x)=1009\textbackslash{}n\textbackslash{}nLet's try to find a pattern. \textbackslash{}nFor x=10, f(x)=1. \textbackslash{}nFor x={-}1000, f(x)=1009. \textbackslash{}nFor x={-}1, f(x)=11. \textbackslash{}nFor x=0, f(x)=9. \textbackslash{}nFor x=1, f(x)=9.\textbackslash{}n\textbackslash{}nLet's look at f(x) + x:\textbackslash{}n0+9=9\textbackslash{}n1+9=10\textbackslash{}n10+1=11\textbackslash{}n{-}1+11=10\textbackslash{}n{-}1000+1009=9\textbackslash{}n\textbackslash{}nThe sum f(x)+x is 9, 10, 11, 10, 9. This looks like it might be related to the number of digits} \\
            }
        }
    }
    & & \\ \\

    \theutterance \stepcounter{utterance}  
    & & & \multicolumn{4}{p{0.6\linewidth}}{
        \cellcolor[rgb]{0.9,0.9,0.9}{
            \makecell[{{p{\linewidth}}}]{
                \texttt{\tiny{[P1$\langle$GM]}}
                \texttt{Game Violation: ParseError: Response is not valid JSON.} \\
\texttt{Output ONLY:} \\
\texttt{One JSON with confidence\_score, current\_mode, input\_rationale, and hypothesis keys ...} \\
\texttt{OR} \\
\texttt{SOLVE: ```python ... ```} \\
            }
        }
    }
     \\ \\

    \theutterance \stepcounter{utterance}  
    & \multicolumn{4}{p{0.6\linewidth}}{
        \cellcolor[rgb]{0.561,0.737,0.561}{
            \makecell[{{p{\linewidth}}}]{
                \texttt{\tiny{[P1$\rangle$GM]}}
                \texttt{\{} \\
\texttt{  "current\_mode": "probing",} \\
\texttt{  "input\_rationale": "I have f(0)=9, f(1)=9, f(10)=1, f({-}1)=11, f({-}1000)=1009. Let's look at the relationship between x and f(x) again. \textbackslash{}n\textbackslash{}nx=0, f(x)=9\textbackslash{}nx=1, f(x)=9\textbackslash{}nx=10, f(x)=1\textbackslash{}nx={-}1, f(x)=11\textbackslash{}nx={-}1000, f(x)=1009\textbackslash{}n\textbackslash{}nLet's try to find a pattern. \textbackslash{}nFor x=10, f(x)=1. \textbackslash{}nFor x={-}1000, f(x)=1009. \textbackslash{}nFor x={-}1, f(x)=11. \textbackslash{}nFor x=0, f(x)=9. \textbackslash{}nFor x=1, f(x)=9.\textbackslash{}n\textbackslash{}nLet's look at f(x) + x:\textbackslash{}n0+9=9\textbackslash{}n1+9=10\textbackslash{}n10+1=11\textbackslash{}n{-}1+11=10\textbackslash{}n{-}1000+1009=9\textbackslash{}n\textbackslash{}nThe sum f(x)+x is 9, 10, 11, 10, 9. This looks like it might be related to the number of digits} \\
            }
        }
    }
    & & \\ \\

    \theutterance \stepcounter{utterance}  
    & & & \multicolumn{2}{p{0.3\linewidth}}{
        \cellcolor[rgb]{0.9,0.9,0.9}{
            \makecell[{{p{\linewidth}}}]{
                \texttt{\tiny{[GM$|$GM]}}
                \texttt{Game Verdict: LOSS (turn limit reached)} \\
            }
        }
    }
    & & \\ \\

\end{supertabular}
}

\end{document}
